\subsection{有界扩散的复合}

命题 13. --- 设 T, X, Y 是三个局部紧空间，\(\Lambda : t \mapsto \lambda_t\) 是 T 在 X 中的一个有界扩散，H: \(x \mapsto \eta_x\) 是 X 在 Y 中的一个有界扩散。则映射 \(t \mapsto \lambda_t\mathrm{H}\) 是 T 在 Y 中的一个有界扩散，记作 \(\Lambda\mathrm{H}\) ，并且

\[
\| \Lambda H \| \leqslant \| \Lambda \| \| H \|.\tag{14}
\]

设 \(f\) 是一个普遍可测函数 \(\geqslant 0\) （定义在 \(Y\) 上），\(\mu\) 是 \(T\) 上一个测度。假设 \(\mu\) 属于 \(\Lambda\) 的定义域，且 \(\mu \Lambda\) 属于 H 的定义域；则 \(\mu\) 属于 \(\Lambda \mathrm{H}\) 的定义域，并且

\[
\begin{array}{l} \langle \mu (\Lambda \mathrm{H}), f \rangle = \langle \mu \Lambda , \mathrm{H} f \rangle = \langle \mu , \Lambda \mathrm{H} f \rangle ; \\ (\mu \Lambda) \mathrm{H} = \mu (\Lambda \mathrm{H}); \quad \Lambda (\mathrm{H} f) = (\Lambda \mathrm{H}) f. \end{array}\tag{15}
\]

置 \(\gamma_{t} = \lambda_{t}H\) ；我们将用 \(\Gamma\) 表示映射 \(\Lambda H\) （T 到 \(\mathcal{M}_{+}(Y)\) 中的），用 \(\Gamma f\) 表示函数 \(t \mapsto \langle \gamma_{t}, f \rangle\) （这是记号的滥用，因为我们尚不知道 \(\Gamma\) 是否为扩散）。由 (13)，\(\langle \gamma_{t}, f \rangle = \langle \lambda_{t}H, f \rangle = \langle \lambda_{t}, Hf \rangle\) ；由于函数 Hf 在 X 上是正的且普遍可测的（命题 12 的推论），首先得出 \(\Gamma f = \Lambda(Hf)\) ，继而得出 \(\Gamma f\) 在 T 上普遍可测（同处引证）。显然，所有测度 \(\gamma_{t}\) 的总质量至多等于 \(\|\Lambda\| \|H\|\) 。因此 \(\Gamma g\) 对每个函数 \(g \in \mathcal{K}_{+}(Y)\) 都是普遍可测且有界的；从而 \(\Gamma\) 对 T 上每个有界测度、特别地对每个具有紧支集的测度，是标量本质可积的。更一般地，如果 \(\mu\) 是 \(\Lambda\) 的定义域中的一个测度，使得 \(\mu\Lambda\) 属于 H 的定义域，那么对 \(g \in \mathcal{K}_{+}(Y)\) ，

\[
\langle \mu , \Gamma g \rangle = \langle \mu , \Lambda (\mathrm{H} g) \rangle = \langle \mu \Lambda , \mathrm{H} g \rangle = \langle (\mu \Lambda) \mathrm{H}, g \rangle .
\]

由于最后的量是有限的，我们看出 \(\Gamma\) 是标量本质 \(\mu\) -可积的。让我们用 \(\mu\Gamma\) 表示积分 \(\int\gamma_{t}d\mu(t)\) （这是记号的滥用，因为我们尚不知道 \(\Gamma\) 是否为扩散）。于是前面的关系式可以写成

\[
\langle \mu \Gamma , g \rangle = \langle (\mu \Lambda) \mathrm{H}, g \rangle ,
\]

或者写成 \(\mu \Gamma = (\mu \Lambda)\mathrm{H}\) ，因为 \(g\) 在 \(\mathcal{K}_+(\mathrm{Y})\) 中是任意的。

再次考虑普遍可测函数 \(f \geqslant 0\) 。我们有

\[
\langle \mu \Gamma , f \rangle = \langle (\mu \Lambda) \mathrm{H}, f \rangle = \langle \mu \Lambda , \mathrm{H} f \rangle = \langle \mu , \Lambda (\mathrm{H} f) \rangle = \langle \mu , \Gamma f \rangle .
\]

当 f 下半连续且 \(\mu\) 跑遍具有紧支集的正测度所成的集时，这些关系式表明 \(\Gamma\) 是 T 在 Y 中的一个扩散。于是该论断不过是把上面证明过程中得到的关系式明显化而已。

定义 4. --- 记号沿用命题 13，扩散 \(\Lambda H\) 称为有界扩散 \(H\) 与 \(\Lambda\) 的复合扩散（或复合）。

设 \(X_{1}, X_{2}, X_{3}, X_{4}\) 是四个局部紧空间，\(\Lambda_{1}, \Lambda_{2}, \Lambda_{3}\) 是三个有界扩散，分别为 \(X_{1}\) 在 \(X_{2}\) 中、\(X_{2}\) 在 \(X_{3}\) 中、\(X_{3}\) 在 \(X_{4}\) 中的扩散。由命题 13 立即得出

\[
(\Lambda_ {1} \Lambda_ {2}) \Lambda_ {3} = \Lambda_ {1} (\Lambda_ {2} \Lambda_ {3}).
\]

因此，对于扩散的复合，我们将使用这些不带括号的记号。

例子. --- 设 u 是 T 到 X 中的一个普遍可测映射，v 是 X 到 Y 中的一个普遍可测映射；由命题 2 b)，用下列公式定义扩散 \(\Lambda\) 和 H：

\[
\lambda_ {t} = \varepsilon_ {u (t)}, \eta_ {x} = \varepsilon_ {v (x)};
\]

则扩散 \(\Gamma = \Lambda H\) 由下式给出：

\[
\gamma_ {t} = \varepsilon_ {(v \circ u) (t)}.
\]

因此要小心注意：扩散的复合次序与通常函数复合的次序相反。

